\documentclass[10pt,a4paper]{amsart}
\usepackage[margin=19mm]{geometry}
\usepackage{amsmath,amssymb,mathtools}
\usepackage[scr=rsfs,cal=euler]{mathalfa}
\usepackage{xfrac}
\usepackage[unicode,hidelinks,bookmarks=false]{hyperref}
\newcommand{\ie}{i.e.\ }
\newcommand{\cf}{cf.\ }
\newcommand{\resp}{respectively\ }
\newcommand{\Subsection}{\subsection}
\providecommand{\nobreakdash}{\nobreak\hbox{-}}
\setlength{\emergencystretch}{2em}
\multlinegap0pt
\usepackage{tikz}
\usetikzlibrary{cd}
\usepackage{amssymb}
\usepackage{xcolor}
\newtheorem{proposition}{Proposition}
\newtheorem{theorem}{Theorem}
\renewcommand{\S}{\mathcal{S}}
\title{\bf \Large Linking PageRank, Time Reversal, and Policy Evaluation}
\author{Konstantin Avrachenkov, Lorenzo Gregoris, Nelly Litvak}
\date{Source: arXiv:2605.00532v2 (CC BY 4.0)}
\begin{document}
\maketitle

\noindent\emph{Source and licence.} \bf \Large Linking PageRank, Time Reversal, and Policy Evaluation. Version arXiv:2605.00532v2 (CC BY 4.0), distributed by arXiv under the Creative Commons Attribution 4.0 International licence, http://creativecommons.org/licenses/by/4.0/, as recorded in the arXiv licence metadata for this exact version and shown on the abstract page https://arxiv.org/abs/2605.00532v2. Full licence notice: \texttt{licenses/arxiv-2605.00532-CC-BY-4.0.txt}.\par\smallskip

\noindent\textit{Editorial scope.} Counted unit: the article's theorem thm:resolvent\_conjugacy (source0.tex lines 565--589). The article's complete original proof of it is delivered with it (source0.tex lines 591--607, one \texttt{proof} environment of 17 lines). No mathematics is written, repaired or re-derived: the statement and the proof are the article's own lines, in the article's own notation, and no retained line is substituted or re-flowed.  Retained uncounted as the source-local context the proof needs: lines 536--546.  Nothing was omitted from any retained span, so the package carries no omission record.  No retained line contains a citation, so the wrapper carries no bibliography and no bibliography entry.  The retained lines contain the article's own diagram code, which the wrapper typesets with the standard tikz packages; no image file is delivered and the package declares no asset dependency.  Editorial formatting only, in addition: an AMS \texttt{amsart} preamble that restates the article's own declarations and macros used by the retained lines, namely the environments \texttt{proposition}, \texttt{theorem} on separate counters, together with the standard packages those lines need.  The retained environments print in this document as Proposition 1, Theorem 1 in order of appearance, whereas the article prints the same items with the numbering of its own full document.  The complete native source of the article as distributed in the arXiv e-print for this exact version is bundled unchanged as \texttt{sources/199572/source0.tex}.
\begin{proposition}
\label{prop:resolvent_adjointness}
The forward and backward resolvents on $L^2(\S,\mu)$
are adjoint:
\[
\langle\mathcal R_\gamma f,g\rangle_\mu
=
\langle f,\mathcal R_\gamma^*g\rangle_\mu,
\qquad f,g\in L^2(\S,\mu).
\]
\end{proposition}

\begin{theorem}
\label{thm:resolvent_conjugacy}
For every $f\in L^2(\S,\mu)$,
\[
\mathcal R_\gamma^*(\varphi f)
=
\varphi(\mathcal R_\gamma f).
\]
Equivalently, the following diagram commutes:
\[
\begin{tikzcd}
L^2(\S,\mu)
    \arrow[r, "\mathcal R_\gamma"]
    \arrow[d, "\varphi"']
&
L^2(\S,\mu)
    \arrow[d, "\varphi"]
\\
\mathcal M(\S,\mu)
    \arrow[r, "\mathcal R_\gamma^*"']
&
\mathcal M(\S,\mu)
\end{tikzcd}
\]
\end{theorem}

\begin{proof}
Let $f,g\in L^2(\S,\mu)$. Using the dual action on
measures, the definition of $\varphi$, and
Proposition~\ref{prop:resolvent_adjointness}, we obtain
\[
\begin{aligned}
\langle g,\mathcal R_\gamma^*(\varphi f)\rangle
&=\langle\mathcal R_\gamma^*g,\varphi f\rangle\\
&=\langle f,\mathcal R_\gamma^*g\rangle_\mu\\
&=\langle\mathcal R_\gamma f,g\rangle_\mu\\
&=\langle g,\varphi(\mathcal R_\gamma f)\rangle.
\end{aligned}
\]
Since this holds for every $g$, taking
$g=\mathbf 1_{\{s\}}$ for each $s\in\S$ shows
that the two measures are equal.
\end{proof}

\end{document}
